\subsection{Evaluasi Kinerja Model dan Sistem}

Evaluasi dilakukan untuk menilai kemampuan YOLO11-seg dalam mengenali objek serta kemampuan sistem dalam mengolah hasil prediksi menjadi jumlah lauk dan total harga. Evaluasi model mencakup hasil deteksi dan segmentasi, sedangkan evaluasi sistem mencakup ketepatan penghitungan instans dan kecepatan pemrosesan.

\subsubsection*{a. Evaluasi Deteksi dan Segmentasi}

Kesesuaian antara prediksi dan \textit{ground truth} dapat diukur menggunakan \textit{Intersection over Union} (IoU). IoU merupakan perbandingan antara luas irisan dan luas gabungan wilayah prediksi dengan \textit{ground truth}, dan digunakan dalam evaluasi model segmentasi berbasis YOLO \cite{qiu2025}. Jika $A$ merupakan wilayah prediksi dan $B$ merupakan wilayah \textit{ground truth}, IoU dinyatakan sebagai:

\begin{equation}
    IoU(A,B) = \frac{|A \cap B|}{|A \cup B|}
    \label{eq:iou}
\end{equation}

Berdasarkan pencocokan antara prediksi dan \textit{ground truth}, hasil prediksi dapat dikelompokkan menjadi \textit{true positive} (TP), \textit{false positive} (FP), dan \textit{false negative} (FN). \textit{Precision}, \textit{recall}, dan \textit{F1-score} merupakan metrik yang umum digunakan untuk mengevaluasi hasil deteksi dan \textit{instance segmentation}. Rumus ketiga metrik tersebut juga digunakan dalam penelitian \textit{instance segmentation} berbasis YOLOv8-Seg \cite{yue2023}.

\textit{Precision} menunjukkan proporsi prediksi positif yang benar dan dinyatakan sebagai:

\begin{equation}
    \textit{Precision} = \frac{TP}{TP + FP}\times 100\%
    \label{eq:precision}
\end{equation}

\textit{Recall} menunjukkan proporsi objek sebenarnya yang berhasil dikenali dan dinyatakan sebagai:

\begin{equation}
    \textit{Recall} = \frac{TP}{TP + FN}\times 100\%
    \label{eq:recall}
\end{equation}

Nilai \textit{F1-score} merupakan rata-rata harmonik dari \textit{precision} dan \textit{recall}, yang dinyatakan sebagai:

\begin{equation}
    F1 = 2 \times
		{\textit{Precision}} \times
    \frac{\textit{Recall}}
         {\textit{Precision} + \textit{Recall}}
    \label{eq:f1}
\end{equation}

Selain itu, performa YOLO11-seg dinilai menggunakan \textit{mean Average Precision} (mAP). mAP50 menunjukkan nilai mAP pada ambang IoU 0,50, sedangkan mAP50--95 merangkum performa pada rentang ambang IoU 0,50 hingga 0,95. Pada model segmentasi, metrik mAP dapat dilaporkan untuk hasil \textit{bounding box} dan \textit{segmentation mask} \cite{qiu2025,syafira2025}.

\subsubsection*{b. Evaluasi Penghitungan Instans}

Metrik deteksi dan segmentasi belum secara langsung menunjukkan seberapa tepat jumlah objek yang dihasilkan sistem. Pada penelitian SibNet, \textit{Mean Absolute Error} (MAE) digunakan untuk mengevaluasi kesalahan penghitungan instans makanan berdasarkan selisih absolut antara jumlah prediksi dan jumlah acuan \cite{nguyen2022sibnet}. Pada penelitian ini, bentuk MAE tersebut disesuaikan untuk menghitung kesalahan penghitungan pada masing-masing kelas lauk. Untuk kelas ke-$i$ pada $N$ citra pengujian, MAE dinyatakan dengan Persamaan~\ref{eq:mae-count}.

\begin{equation}
    MAE_i = \frac{1}{N} \sum_{j=1}^{N} \left|\hat{n}_{ji} - n_{ji}\right|
    \label{eq:mae-count}
\end{equation}

Pada persamaan tersebut, $i$ menunjukkan kelas yang sedang dinilai, $j$ menunjukkan citra pengujian, $\hat{n}_{ji}$ merupakan jumlah instans hasil prediksi pada citra ke-$j$ untuk kelas ke-$i$, sedangkan $n_{ji}$ merupakan jumlah instans berdasarkan \textit{ground truth}. Nilai MAE yang lebih kecil menunjukkan rata-rata kesalahan penghitungan yang lebih rendah, sedangkan nilai nol menunjukkan tidak terdapat selisih jumlah pada data yang dievaluasi.

\subsubsection*{c. Kecepatan Inferensi}

Kecepatan model dapat dinilai berdasarkan waktu inferensi rata-rata yang diperlukan untuk memproses satu citra. Penggunaan waktu inferensi sebagai salah satu indikator performa juga ditemukan pada penelitian segmentasi berbasis YOLO untuk menilai kesesuaian model terhadap kebutuhan pemrosesan secara \textit{real-time} \cite{syafira2025}. Perangkat, data pengujian, dan prosedur pengukuran yang digunakan dijelaskan pada Bab III.
